\section{Technical Concept}
\label{sec:technical_concept}

The technical concept of BiteBound is designed to ensure real-time responsiveness, modularity, and thread safety. By leveraging the dual-core capabilities of the ESP32-S3, the system separates latency-sensitive gameplay and rendering tasks from time-intensive network operations.

\subsection{High-Level Architecture}
\label{subsec:high_level_architecture}

The architecture relies on an asynchronous, event-driven design. The ESP32-S3 runs two primary execution threads across its dual cores to balance the computational workload. Core 1 is designated as the Game Core, measuring the sensor data, executing the physics engine, handling the game logic, and updating the display. Core 0 acts as the Network Core, managing secure TLS connections, MQTT messaging, and transmitting telemetry, running significantly slower than the Game Core.

Communication between the embedded hardware and the external control interfaces is bridged through the centralized MQTT broker HiveMQ~\cite{hivemq2026,mqtt2026}. External clients communicate bidirectionally with the device. They send game-related commands to a dedicated topic, while subscribing to a telemetry topic to monitor real-time variables.

The system architecture is depicted in Figure~\ref{fig:system-architecture}, illustrating the interactions between the ESP32-S3 and the frontend dashboards (see \texttt{diagrams/out/}).

\begin{figure}[h]
    \centering
    \includegraphics[width=0.9\linewidth]{../diagrams/out/system-architecture/system-architecture.png}
    \caption[BiteBound Architecture]{System architecture diagram displaying the user interactions with the ESP32-S3 and the frontend dashboards. }
    \label{fig:system-architecture}
\end{figure}

\subsection{Components}
\label{subsec:components_concept}

The system architecture is divided into three primary components:

\begin{itemize}
    \item \textbf{Hardware (ESP32-S3):} The core of the system is the Waveshare ESP32-S3 development board equipped with a 1.69-inch touch LCD, a 6-axis inertial measurement unit (IMU), and a buzzer. It hosts the procedural sensor manager, maze generator, 2D physics simulation, game engine, display logic, and network communication stack.
    \item \textbf{Android Application:} A native mobile interface written in Kotlin. It receives telemetry data and sends control commands to the ESP32-S3 via MQTT.
    \item \textbf{Node-RED Dashboard:} A flow-based UI-Builder dashboard that mirrors the application logic of the Android app.
\end{itemize}